\begingroup
\setlength{\LTleft}{0pt}
\setlength{\LTright}{0pt}
\renewcommand{\arraystretch}{1.3}
\begin{longtable}{@{}p{0.20\textwidth} p{0.385\textwidth} p{0.385\textwidth}@{}}
\caption[Resumen de las características del Ag Barometer Austral y del Ag Economy Barometer]{\label{tbl-comparacion}Resumen de las características del \textit{Ag Barometer Austral} y del \textit{Ag Economy Barometer}}\\
\toprule
\textbf{Característica} & \centering\arraybackslash\textit{\textbf{Ag Barometer Austral}} & \centering\arraybackslash\textit{\textbf{Ag Economy Barometer}} \\
\midrule
\endfirsthead
\multicolumn{3}{@{}l}{\small (continuación)} \\
\toprule
\textbf{Característica} & \centering\arraybackslash\textit{\textbf{Ag Barometer Austral}} & \centering\arraybackslash\textit{\textbf{Ag Economy Barometer}} \\
\midrule
\endhead

\bottomrule
\multicolumn{3}{@{}p{\textwidth}@{}}{\footnotesize\raggedright Fuente: Elaboración propia en base a datos de la Universidad Austral y Universidad Purdue} \\
\endlastfoot

\textbf{Población} &
Conformada por productores:
\newline\textbullet\ De la Pampa Húmeda
\newline\textbullet\ Que disponen de un valor anual de producción mayor o igual a USD 200.000
\newline\textbullet\ Donde un 45\% producen soja, un 21\% maíz, un 16\% ganadería vacuna, un 7\% trigo, un 7\% tambo y un 4\% girasol &
Conformada por productores:
\newline\textbullet\ Que disponen de un valor anual de producción mayor o igual a USD 500.000
\newline\textbullet\ Donde un 53\% producen soja, un 14\% maíz, un 3\% algodón, un 19\% ganadería vacuna, un 5\% tambo y un 6\% ganadería porcina \\[0.5em]

\cmidrule{1-3}
\textbf{Muestreo} &
Se aplicó un muestreo estratificado según la producción anual:
\newline\textbullet\ 50\% de los encuestados producen entre USD 200.000 y USD 399.000
\newline\textbullet\ 30\% producen entre USD 414.000 y USD 999.999
\newline\textbullet\ 20\% producen USD 1.000.000 o más &
Se aplicó un muestreo estratificado según la producción anual:
\newline\textbullet\ 49\% de los encuestados producen entre USD 500.000 y USD 999.999
\newline\textbullet\ 36\% producen entre USD 1.000.000 y USD 2.499.000
\newline\textbullet\ 15\% producen USD 2.500.000 o más \\[0.5em]

\cmidrule{1-3}
\textbf{Trabajo de campo} &
\textbullet\ Se comenzó a realizar en octubre 2018
\newline\textbullet\ El relevamiento se realiza en los 6 meses impares
\newline\textbullet\ Periodicidad bimestral
\newline\textbullet\ La encuesta se realiza por teléfono a 406 productores
\newline\textbullet\ Cada recolección de respuestas insume alrededor de tres semanas &
\textbullet\ Se comenzó a realizar en octubre 2015
\newline\textbullet\ El relevamiento se realiza en los 12 meses
\newline\textbullet\ Periodicidad mensual
\newline\textbullet\ La encuesta se realiza por teléfono a 400 productores
\newline\textbullet\ Cada recolección de respuestas insume aproximadamente una semana \\[0.5em]

\cmidrule{1-3}
\textbf{Cuestionario} &
\textbullet\ Cinco preguntas fijas
\newline\textbullet\ Aproximadamente seis preguntas que recogen respuestas sobre temas coyunturales (no se las utiliza en la construcción de los índices) &
\textbullet\ Cinco preguntas fijas \\[0.5em]

\cmidrule{1-3}
\textbf{Construcción de índices} &
\textbullet\ Primer paso: calcular una medida para cada una de las preguntas.
\newline\textbullet\ Segundo paso: \textit{Ag Barometer Austral} es el promedio simple de las cinco medidas.
\newline\textbullet\ Siguiendo los mismos pasos también se construyen los índices ``Condiciones presentes'' y ``Expectativas futuras'' &
\textbullet\ Primer paso: calcular una medida para cada una de las preguntas.
\newline\textbullet\ Segundo paso: \textit{Ag Economy Barometer} es la suma de las cinco medidas dividida por el período base.
\newline\textbullet\ Siguiendo los mismos pasos también se construyen los índices ``Condiciones presentes'', ``Expectativas futuras'' e ``Inversión del capital agrícola'' \\[0.5em]

\cmidrule{1-3}
\textbf{Interpretación} &
\textbullet\ \textit{Ag Barometer Austral} = 100, valor neutro
\newline\textbullet\ \textit{Ag Barometer Austral} > 100, percepciones optimistas
\newline\textbullet\ \textit{Ag Barometer Austral} < 100, percepciones pesimistas &
\textbullet\ \textit{Ag Economy Barometer} = 100, no hubo cambios en las percepciones respecto del período base
\newline\textbullet\ \textit{Ag Economy Barometer} > 100, aumento en las percepciones respecto del período base
\newline\textbullet\ \textit{Ag Economy Barometer} < 100, disminución en las percepciones respecto del período base \\

\end{longtable}


\endgroup